\documentclass[aspectratio=169,11pt]{beamer}

\usetheme{Madrid}
\usecolortheme{default}
\usefonttheme{professionalfonts}
\setbeamertemplate{navigation symbols}{}
\setbeamertemplate{footline}[frame number]

\usepackage{amsmath, amssymb, booktabs, graphicx}

\title{Institutional Reform, Implementation, and Labor-Market Adjustment}
\subtitle{Labor Markets and Political/Cultural Institutions -- Week 8}
\author{Teaching deck draft}
\date{}

\begin{document}

\begin{frame}
  \titlepage
\end{frame}

\begin{frame}{Central question}
\begin{itemize}
    \item Why do similar institutional reforms produce different labor-market effects across places, firms, and workers?
    \item Which part of the reform chain is actually moving: law, implementation, knowledge, compliance, or equilibrium adjustment?
    \item What labor margin is observed: wages, hiring, firing, contract form, formality, promotion, or mobility?
\end{itemize}
\end{frame}

\begin{frame}{Distinction from Week 7}
\begin{itemize}
    \item Week 7: why old institutions persist through skills, mobility, networks, public goods, and employer power.
    \item Week 8: how new reforms work when implementation and adjustment are uneven.
    \item The object is not legal text alone.
    \item The object is the chain from reform to worker and firm behavior.
\end{itemize}
\end{frame}

\begin{frame}{Reform taxonomy}
\small
\begin{tabular}{p{0.30\linewidth} p{0.30\linewidth} p{0.30\linewidth}}
\toprule
Reform family & Typical labor margin & Why effects differ \\
\midrule
Dismissal protection & firing, hiring, turnover & coverage, legal salience, firm size \\
Enforcement expansion & compliance, formality, wages & local capacity, inspections, evasion \\
Legal information & take-up, claims, bargaining & baseline ignorance, trust, retaliation risk \\
Formalization reform & registration, payroll, contracts & private benefits, taxes, enforcement threat \\
Equal opportunity & hiring, promotion, gaps & discretion, procedures, outside options \\
Transparency & wages, bargaining, pay gaps & disclosure scope, comparison, employer response \\
Administrative reform & matching, benefits, job-finding & agency quality, rollout, employer engagement \\
\bottomrule
\end{tabular}
\end{frame}

\begin{frame}{Transmission framework}
\small
\[
\begin{aligned}
\Delta Y^{labor}
= f(&\text{legal text},\ \text{implementation},\ \text{knowledge},\ \text{exposure},\\
&\text{firm response},\ \text{worker response},\ \text{equilibrium})
\end{aligned}
\]
\begin{itemize}
    \item Legal text changes rights, costs, procedures, disclosure, or administration.
    \item Implementation determines whether the change becomes credible.
    \item Workers and firms respond only if they are exposed and understand the margin.
    \item Equilibrium adjustment changes incidence and welfare.
\end{itemize}
\end{frame}

\begin{frame}{Where heterogeneity enters}
\small
\begin{tabular}{p{0.25\linewidth} p{0.33\linewidth} p{0.32\linewidth}}
\toprule
Stage & Core question & Heterogeneity \\
\midrule
Legal change & What changed on paper? & scope, exemptions, rollout timing \\
Implementation & Was it enforced or administered? & capacity, inspectors, courts, agency quality \\
Knowledge & Did actors understand it? & legal awareness, literacy, trust \\
Exposure & Who is covered? & firm size, sector, contract type, worker type \\
Firm response & How do firms adjust? & compliance, evasion, outsourcing, wage compression \\
Worker response & How do workers react? & search, claims, mobility, bargaining \\
Equilibrium & What else moves? & spillovers, prices, composition, inequality \\
\bottomrule
\end{tabular}
\end{frame}

\begin{frame}{Employment protection and dismissal reform}
\begin{itemize}
    \item Wrongful-discharge doctrines change expected dismissal cost and litigation risk.
    \item Autor--Donohue--Schwab: staggered state adoption and employment consequences.
    \item Observed margins: employment, hiring and firing adjustment, turnover, contract composition.
    \item Interpretation depends on firm size, coverage, worker type, product-market pressure, and legal salience.
\end{itemize}
\end{frame}

\begin{frame}{Implementation and enforcement}
\begin{itemize}
    \item Enforcement expansions change the probability that noncompliance is detected or sanctioned.
    \item Almeida--Carneiro: enforcement of labor regulation and informality in Brazil.
    \item Observed margins: formality, compliance, wages, employment.
    \item The same enforcement increase can bind large formal firms and miss small informal employers.
\end{itemize}
\end{frame}

\begin{frame}{Worker beliefs and legal knowledge}
\begin{itemize}
    \item Rights do not bind if workers do not know them or cannot invoke them.
    \item Bertrand--Crepon: randomized labor-law information in South Africa.
    \item Observed margins: knowledge, beliefs, take-up, and labor-market behavior.
    \item Firm response matters: compliance, screening, postings, contract choice, or avoidance.
\end{itemize}
\end{frame}

\begin{frame}{Firm adaptation and compliance}
\begin{itemize}
    \item Firms can comply, evade, outsource, automate, relabel contracts, or change hiring.
    \item Compliance costs can land on profits, prices, wages, employment, entry, or exit.
    \item Adaptation is heterogeneous by size, visibility, sector, and bargaining environment.
    \item The observed firm margin often determines the welfare interpretation.
\end{itemize}
\end{frame}

\begin{frame}{Formalization and informality margins}
\begin{itemize}
    \item Formalization reforms change registration, payroll reporting, inspections, and benefits access.
    \item Ulyssea: informality as multiple firm and worker margins, not one informal sector.
    \item de Andrade--Bruhn--McKenzie: information, assistance, and enforcement pressure in formalization.
    \item Registration may not imply formal jobs or stronger worker protection.
\end{itemize}
\end{frame}

\begin{frame}{Equal opportunity and procedural reform}
\begin{itemize}
    \item Procedural reforms change who is observed, evaluated, shortlisted, hired, or promoted.
    \item Goldin--Rouse: blind auditions and hiring procedure.
    \item The labor margin is advancement through an evaluation process.
    \item Effects depend on evaluator discretion, occupation structure, and outside options.
\end{itemize}
\end{frame}

\begin{frame}{Transparency reforms}
\begin{itemize}
    \item Transparency changes the information environment for workers and firms.
    \item Baker--Halberstam--Kroft--Mas--Messacar: pay transparency and gender gaps.
    \item Blundell--Duchini--Simion--Turrell: transparency, equality, and pay adjustment.
    \item Margins include wage gaps, pay compression, bargaining, promotions, quits, and firm response.
\end{itemize}
\end{frame}

\begin{frame}{Administrative reform}
\begin{itemize}
    \item Some reforms change public administration rather than private contracting rules.
    \item Public employment services, benefit administration, matching systems, and decentralized bargaining all matter.
    \item Dustmann--Fitzenberger--Schonberg--Spitz-Oener: Germany as broader institutional adjustment.
    \item Main caution: reform packages are bundled, so the operative margin must be named.
\end{itemize}
\end{frame}

\begin{frame}{Empirical design map}
\small
\begin{tabular}{p{0.31\linewidth} p{0.30\linewidth} p{0.29\linewidth}}
\toprule
Setting & Typical design & Observed margin \\
\midrule
State labor-law adoption & staggered DID / event study & employment, turnover, contract choice \\
Enforcement expansion & local intensity variation & compliance, formality, wages \\
Legal information & randomized intervention & knowledge, take-up, behavior \\
Formalization program & field experiment or rollout & registration, payroll, contracts \\
Hiring procedure & event study or experiment & shortlist, hiring, promotion \\
Pay transparency & threshold DID / distributional analysis & gaps, bargaining, compression \\
Administrative reform & phased rollout / DID & matching, job-finding, wages \\
\bottomrule
\end{tabular}
\end{frame}

\begin{frame}{Welfare and incidence}
\begin{itemize}
    \item Take-up: do workers and firms use the reform?
    \item Compliance: does the rule change behavior or only paperwork?
    \item Incidence: who pays through wages, profits, prices, employment, or entry?
    \item Spillovers: what happens to uncovered workers, informal firms, and local markets?
    \item Inequality: does the reform change gaps, ranks, or access to protection?
\end{itemize}
\end{frame}

\begin{frame}{Research frontier}
\begin{itemize}
    \item Link legal text, administrative capacity, worker beliefs, and firm-worker data.
    \item Measure adaptation margins rather than only mean employment effects.
    \item Study reforms in markets with informality, monopsony, subcontracting, and weak courts.
    \item Estimate spillovers to uncovered firms and workers.
    \item Build welfare measures that include information, protection, inequality, and displacement.
\end{itemize}
\end{frame}

\begin{frame}{Week 8 lab}
\begin{itemize}
    \item Reproduce: synthetic legal-information intervention inspired by Bertrand--Crepon.
    \item Diagnose: classify bottlenecks in the reform-transmission chain.
    \item Transfer: map designs to dismissal protection, enforcement, formalization, transparency, and administration.
    \item Output: treatment table, bottleneck diagnostic, transfer design classification.
\end{itemize}
\end{frame}

\begin{frame}{Bridge to Week 9}
\begin{itemize}
    \item Week 8 explains why implementation and adjustment make reform effects heterogeneous.
    \item Week 9 asks which lessons travel across development settings, welfare states, informal economies, and global production systems.
    \item Comparative institutions require portability questions: what changes when capacity, informality, culture, hierarchy, and market structure differ?
    \item The next step is to compare institutional regimes, not only individual reforms.
\end{itemize}
\end{frame}

\begin{frame}{Takeaways}
\begin{itemize}
    \item Reform effects are produced by transmission chains, not legal text alone.
    \item Worker knowledge, firm adaptation, local capacity, and coverage determine what binds.
    \item Empirical designs identify specific margins of reform exposure.
    \item Welfare requires take-up, compliance, incidence, spillovers, inequality, and unintended margins.
\end{itemize}
\end{frame}

\end{document}
